\section{The reduced five-row certificate}\label{app:certificate}

The even-column state orbits are listed in
\eqref{eq:even-orbits}. The odd-column orbits, in order, are
\[
 \{1,16\},\ \{4\},\ \{7,28\},\ \{13,22\},\
 \{19,25\},\ \{31\}.
\]
These twenty states form the set reachable from the empty
state. This can be checked directly from \eqref{eq:Tm}:
starting at mask zero and listing all successors reaches
exactly these states, and all their successors remain in
the set. Row reversal commutes with the transition relation
and fixes the initial state, so reversed states have equal
counts.

Let $u_j$ be the even-column count vector, with one-state
counts as in Section \ref{sec:five}, and let $z_j$ be the next
odd-column vector. Then $z_j=Bu_j$ and $u_{j+1}=Cz_j$, where
\begin{equation}\label{eq:BC}
\begin{split}
 B&=\begin{pmatrix}
1&1&1&1&1&0\\
1&2&0&0&0&1\\
1&1&0&0&0&0\\
0&0&0&1&0&0\\
1&0&1&0&0&0\\
1&0&0&0&0&0
\end{pmatrix},\qquad
 C=\begin{pmatrix}
2&1&2&0&2&1\\
1&1&1&0&0&0\\
1&0&0&0&1&0\\
1&0&0&1&0&0\\
1&0&0&0&0&0\\
0&1&0&0&0&0
\end{pmatrix}.
\end{split}
\end{equation}
For an input orbit $O_j$ and an output representative $q_i$,
the corresponding entry is
\[
\sum_{p\in O_j}T_5[p,q_i].
\]
It does not depend on the choice of representative.
For example, $B_{2,2}=2$ because both masks 3 and 24
can precede mask 4.
The reduced matrices need not be symmetric:
their variables record the count of each state, not the
total of an orbit. Matrix multiplication gives $CB=M$.

The even orbit-size vector is $(1,2,2,2,2,1)$ and the odd
orbit-size vector is $(2,1,2,2,2,1)$.
The fixed masks of row reversal are 0 and 27 on the even side,
and 4 and 31 on the odd side. Consequently,
\[
 w_{v,o}=(2,1,2,2,2,1)B,\qquad
 w_{r,o}=(0,1,0,0,0,1)B,
\]
which gives the terminal vectors in the text.

To check \eqref{eq:certificate} coefficientwise, write
$N(t)=n_0+n_1t+n_2t^2+n_3t^3$, with
\[
\begin{array}{ll}
n_0=(1,0,0,0,0,0)^{\trans},&
n_1=(-7,3,2,1,1,1)^{\trans},\\
n_2=(7,-4,-5,1,0,0)^{\trans},&
n_3=(-1,1,1,0,-1,-1)^{\trans}.
\end{array}
\]
Direct integer multiplication gives
\[
\begin{aligned}
n_1-Mn_0&=-15\evec,&
n_2-Mn_1&=32\evec,\\
n_3-Mn_2&=-15\evec,&
-Mn_3&=\evec.
\end{aligned}
\]
Together with $n_0=\evec$, these are all coefficients of
$(I-tM)N(t)=P(t)\evec$.
Thus the certificate consists of finitely many explicit
integer identities and determines every coefficient of
the generating functions.
